% Software architecture (input file of touchsynth_poly_report.tex). Draft for review, 2026-10-20.

% =====================================================================================================================
\section{Software architecture}\label{sec:sad}

This section describes how the software is built to meet the requirements of Section~\ref{sec:srs}. It covers where
each part runs, how it is divided into modules, what runs in which execution context with which budget, how data flows
from a MIDI event to the DAC, and how the software is configured and verified. Section~\ref{sec:fw} gives the
details of the firmware as built; the voice bus is specified in Section~\ref{sec:bus}.

\review{First draft. Read the card firmware views as a description of working code (measured on the bench). Read the
controller views as a target design: only the frame engine, the bus scheduler, USB-MIDI and CDC, and the SWIO bridge
exist (\file{firmware/bench\_ctrl}). The decisions in Table~\ref{tab:sadadr} and the questions in
Section~\ref{sec:sadopen} are the points to confirm.}

% ---------------------------------------------------------------------------------------------------------------------
\subsection{Context and deployment}

\begin{figure}[htbp]
\centering
\resizebox{\textwidth}{!}{%
\begin{tikzpicture}[font=\small\sffamily, >={Stealth[length=2mm]},
  blk/.style={draw, rounded corners=2pt, fill=#1, minimum height=9mm, align=center}, blk/.default=white,
  lab/.style={font=\scriptsize}]
  \node[blk=gray!6, minimum width=38mm] (link) at (0,3.9) {WCH-LinkE\\\scriptsize development only};
  \node[blk=gray!12, minimum width=38mm, minimum height=24mm] (pc) at (0,0.9)
    {\textbf{computer}\\[2pt]\scriptsize DAW, keyboard software\\\scriptsize editor / librarian (CDC, SysEx)\\
     \scriptsize minichlink, test suite};
  \node[blk=gray!12, minimum width=38mm] (kbd) at (0,-2.3) {USB-MIDI keyboard\\\scriptsize (USB-A host port)};
  \node[blk=blue!18, minimum width=50mm, minimum height=50mm] (ct) at (7.2,0.4)
    {\textbf{controller firmware}\\\scriptsize CH32V203, 144\,MHz, 64\,KB / 20\,KB\\[4pt]
     \scriptsize USB device (MIDI + CDC), USB host\\\scriptsize MIDI router, parts, allocator\\
     \scriptsize bus scheduler, frame engine\\\scriptsize enumeration, supervision\\\scriptsize patch store,
     calibration store\\\scriptsize SWIO programmer};
  \foreach \n/\y in {1/2.6, 2/1.1, 8/-2.2}
    \node[blk=yellow!20, minimum width=48mm, minimum height=11mm] (c\n) at (15.6,\y)
      {\textbf{card firmware} (slot \n)\\\scriptsize CH32V006, 48\,MHz HSI, 62\,KB / 8\,KB};
  \node at (15.6,-0.45) {\large$\vdots$};
  \draw[<->] (1.9,0.9) -- node[above,lab]{USB-MIDI, CDC} (4.7,0.9);
  \draw[->] (kbd.east) -- node[above,lab]{USB host} (kbd.east -| ct.west);
  \draw[->, gray] (link.east) -- node[above,lab]{SWD} (link.east -| ct.west);
  \draw[->, thick] (9.7,2.0) -- node[above,lab]{BCAST, SYNC} (12.3,2.0);
  \draw[thick] (12.3,2.6) -- (12.3,-2.2);
  \foreach \n in {1,2,8} \draw[->, thick] (c\n.west -| 12.3,0) -- (c\n.west);
  \fill (12.3,2.0) circle (0.6mm);
  \draw[dashed] (12.8,2.6) -- (12.8,-2.2);
  \foreach \n in {1,2,8} \fill (c\n.west -| 12.8,0) circle (0.5mm);
  \draw[->, dashed] (12.8,0.3) -- node[below,lab]{RETURN} (9.7,0.3);
  \draw[<->] (9.7,-1.8) -- (11.0,-1.8) -- (11.0,-3.3) -- node[below,lab]{NRST$_k$ (74HC595) and SWIO$_k$: one pair
    per slot} (15.6,-3.3) -- (c8.south);
\end{tikzpicture}}
\caption{Software context. Two firmwares on nine microcontrollers. The computer plays and edits the instrument
over USB, and also runs the test suite and the programming tools during development. The controller is the only bus
master and the only owner of the instrument's state; each card runs one voice.}
\label{fig:sadctx}
\end{figure}

\begin{itemize}
  \item \textbf{Card firmware:} one image, identical on every card. A card knows nothing about its slot until it is
    enumerated, and it stores nothing except its calibration. It can be programmed by a WCH-LinkE (development, a card
    on its own) or by the controller over its SWIO line.
  \item \textbf{Controller firmware:} one image on the backplane. It holds the patches, the setups, the part
    configuration and the calibration copies, and is the time reference (crystal, SYNC). It is programmed over SWD
    or by the factory bootloader (BOOT0).
  \item \textbf{Host side:} no driver (class-compliant USB-MIDI, CDC ACM). The editor and librarian are not part
    of this project; CDC and SysEx are their interface.
\end{itemize}

% ---------------------------------------------------------------------------------------------------------------------
\subsection{Architectural drivers}

\begin{enumerate}
  \item \textbf{Alignment over latency} (P8, P10): every card must sound the same event at the same carrier period,
    independent of when its bytes arrive or what it is computing. A constant latency of a few milliseconds is
    acceptable.
  \item \textbf{One voice in \SI{48}{MHz} and \SI{8}{KB}} (C3): the DAC renderer needs one sample per
    \SI{5.33}{\micro s}, the control tick 1\,kHz, the bus parser up to 100 bytes per frame, all without a missed buffer
    half.
  \item \textbf{Robustness without acknowledgements:} lost packets, card resets, controller resets and missing SYNC
    edges must heal by themselves.
  \item \textbf{Testability on the bench:} one real card plus virtual voices must exercise the full protocol and load;
    timing must be observable through counters.
  \item \textbf{Updatable in the field:} cards from the controller, the controller from the computer.
\end{enumerate}

% ---------------------------------------------------------------------------------------------------------------------
\subsection{Principal decisions}

\begingroup\small\setlength{\tabcolsep}{4pt}
\begin{longtable}{@{}>{\raggedright}p{0.7cm} >{\raggedright}p{4.4cm} >{\raggedright}p{3.6cm}
  >{\raggedright\arraybackslash}p{6.9cm}@{}}
\caption{Architecture decisions.}\label{tab:sadadr}\\
\toprule
ID & Decision & Alternative & Reason, consequence \\
\midrule\endfirsthead
\toprule
ID & Decision & Alternative & Reason, consequence \\
\midrule\endhead
\bottomrule\endfoot
A1 & Time-triggered bus: events carry their execution frame ($a + L$); outputs follow $D$ frames later & events run on
  arrival & Alignment to one carrier period on every card, whatever the parse and render timing. Costs a constant
  5--7\,ms latency (P8 allows 10). \\
A2 & The controller owns all state; cards are stateless except for their calibration; periodic refreshes repair
  losses & an acknowledged protocol; patches stored on the cards & Broadcast stays simple, one card speaks at a time,
  a card can be swapped or reset at any time. Costs about 1.5 bytes per frame of refresh. \\
A3 & DAC render in the DMA half / complete interrupt at the low preemption level, with nesting; everything else at
  the high level & render in the main loop & A stall in the main loop (console) can never cost audio, and the tick
  enters within 27--28 cycles during a render. The render's budget must include the ticks that preempt it. \\
A4 & Bus parsing in the control tick, with early discard of foreign packets & parsing in the main loop below the
  render & Keeps $L = 2$. Affordable since the render rewrite (at least 99 periods of slack). The fallback, if card CPU
  becomes tight again, needs $L = 3$. \\
A5 & Bare metal on ch32fun: static buffers, no RTOS, no heap & FreeRTOS or similar & \SI{8}{KB} RAM, four execution
  contexts, deterministic timing; the stack is monitored instead (NF3). \\
A6 & One card code base for board~1 and the voice card (\texttt{BOARD}) & a separate card firmware & The engine,
  DSP and drivers stay tested on both; costs conditional code in \file{hw.c} and \file{board.h}. \\
A7 & The renderer is bit-exact against a reference revision; its hot loops are inline assembly, run from SRAM, and
  verified by an on-target hash & plain C; a half-rate render & Four to five times faster with an unchanged sound. The
  assembly is verified only by the hash, and the RAM code costs \SI{2.3}{KB}. \\
A8 & Controller USB: an own USBD driver and ch32fun's \file{fsusb.c} for USBFS behind one API, flow control by NAK &
  ch32fun \file{usbd.c}; TinyUSB & No message is lost, and DMA1 stays on (\file{usbd.c} switches it off). TinyUSB is
  reconsidered for the host port. \\
A9 & Cards are updated over SWIO from the controller (WCH one-wire debug protocol, bit-banged) & a bootloader on the
  card, over the bus & No card bootloader to maintain or to brick, and the same path as the WCH-LinkE (minichlink).
  Costs 8 controller GPIOs. \\
A10 & NRST through the 74HC595 without pull-downs; cards run freely while the controller is not running & cards held
  in reset & A controller reset or a debugger halt needs no special handling. The firmware must take the cards back
  (reset all, enumerate). \\
A11 & Each card trims its HSI to SYNC and scales its pitch by the measured frame period & a shared \SI{24}{MHz} clock &
  No clock distribution. The HSI's short-term wander (0.35\,cent rms) remains; CLK\_IN stays as an option. \\
A12 & The host software is a browser app (Web MIDI, Web Serial), static files without a build step & a Python
  desktop app; a DAW plugin (VST3, CLAP) & No installation on any system, works offline; the Python tools stay for
  automated tests. Needs a Chromium-based browser (Web Serial). \\
A13 & One parameter table (JSON) generates the app's parameter list, the controller firmware's tables and the docs'
  table & hand-kept copies & The app, the firmware and midi.md cannot drift apart. \\
\end{longtable}
\endgroup

% ---------------------------------------------------------------------------------------------------------------------
\subsection{Card firmware}\label{sec:sadcard}

\paragraph{Modules.} The card firmware has four layers. Each module calls only those below it, except for the
callbacks named in Table~\ref{tab:sadmod}.

\begin{table}[htbp]
\centering\small
\caption{Card firmware modules and their interfaces (\file{firmware/voice/}).}\label{tab:sadmod}
\begin{tabularx}{\textwidth}{@{}l l X@{}}
\toprule
Layer & Module & Responsibility; interface \\
\midrule
application & \file{touchsynth.c} & start-up, interrupt levels, the control tick dispatcher (\texttt{control\_tick}:
  bus, synth, CV commit, $V_n$ commit), SysTick fallback, main loop; stack monitor \\
 & \file{console.c} & SWIO console: diagnostics, manual control, render bench and check \\
protocol & \file{bus.c} & bus state machine (Figure~\ref{fig:sadstate}), HSI trim, frame counter, $\bar P$, in-place
  COBS / CRC parser, event queue, replies; \texttt{bus\_tick}, \texttt{bus\_present} \\
voice & \file{synth.c} & voice engine at \SI{1}{kHz} (envelopes, LFO, glide, CV laws), bus event handlers
  \texttt{synth\_bus\_*}, presets and tunes \\
 & \file{dsp.c} & DAC engines (oscillators, noise, modal), error feedback, $V_n$ queue and block split;
  \texttt{dsp\_pitch}, \texttt{dsp\_strike}, \texttt{dsp\_commit}, the render callback \texttt{dsp\_fill} \\
 & \file{tables.h} & generated: presets, CV laws, mode tables, exponential tables \\
drivers & \file{hw.c} & TIM1 CV PWM and its DMA burst, CV delay line; TIM2 DAC, ring and DMA ch1 with the render
  interrupt; SYNC EXTI; USART1 RX ring and TX by DMA; ADC capture; LEDs; chip ID \\
 & \file{board.h} & pin map and timer remaps per board and chip \\
platform & ch32fun & start-up code, register definitions, \texttt{funPinMode} \\
\bottomrule
\end{tabularx}
\end{table}

\paragraph{Execution contexts.} Table~\ref{tab:sadexec} lists every context with its trigger, level and measured
cost. The budget rule follows from A3: a buffer half's render time plus the ticks that preempt it must stay below
65\,536 cycles. The slack, in carrier periods, is what the counters report.

\begin{table}[htbp]
\centering\small
\caption{Card execution contexts (V006 at \SI{48}{MHz}; costs from the bench, Section~\ref{sec:cpu}).}\label{tab:sadexec}
\setlength{\tabcolsep}{4pt}
\begin{tabularx}{\textwidth}{@{}>{\raggedright}p{1.9cm} >{\raggedright}p{2.9cm} >{\raggedright}p{1.3cm} >{\raggedright}X >{\raggedright\arraybackslash}p{3.7cm}@{}}
\toprule
Context & Trigger, rate & Level & Work & Worst measured \\
\midrule
control tick & SYNC edge, EXTI4 (PD4), \SI{1}{kHz} & high & parse bytes received before the tick, run due
  events, replies; synth tick; CV commit; $V_n$ commit & 7\,k cycles idle, 9\,k at 71\,B/frame, 19\,k with a reply,
  27\,k at a modal PATCH\_SELECT \\
SysTick & only without SYNC (standalone, a missing edge) & high & the same tick & as above \\
DAC render & DMA1 ch1 half / complete, \SI{732}{Hz} & low & render the half that has just played: 256 periods in
  blocks of 32, split at $V_n$ boundaries & 8--12\,k (oscillators), 21\,k (modal); with preempting ticks 25\,k of
  65.5\,k \\
CV dither & TIM1 CC1 request at each update, \SI{46.875}{kHz} & DMA ch2 & 4-register burst from a 16-halfword table &
  no CPU \\
bus RX, TX & USART1 & DMA ch5, ch4 & receive ring of 256 bytes; replies & no CPU \\
main loop & -- & -- & SWIO console, stack low-water mark & -- \\
\bottomrule
\end{tabularx}
\end{table}

\paragraph{Data flow and the presentation delay.} Figure~\ref{fig:sadflow} follows one NOTE\_ON. It arrives during
frame $f-1$, the tick of frame $f$ parses it (its execution frame), and the voice engine computes the new state.
That state is delayed to frame $f + D$: the CVs through a delay line of $D + 1$ entries, the DAC state $V_f$ through a
queue stamped with $b_f = t_f + D\,\bar P$. The renderer splits its block at the carrier period of $b_f$, so the
change is exact to one period on every card.

\begin{figure}[htbp]
\centering
\resizebox{\textwidth}{!}{%
\begin{tikzpicture}[font=\small\sffamily, >={Stealth[length=2mm]},
  st/.style={draw, rounded corners=2pt, fill=#1, minimum height=10mm, minimum width=24mm, align=center},
  st/.default=white, buf/.style={draw, fill=gray!10, minimum height=9mm, minimum width=22mm, align=center},
  lab/.style={font=\scriptsize}]
  \node[buf] (ring) at (0,0) {RX ring\\\scriptsize 256 B, DMA ch5};
  \node[st=yellow!20] (parse) at (3.3,0) {parser\\\scriptsize COBS, CRC, dst};
  \node[buf] (evq) at (6.6,0) {event queue\\\scriptsize 16 events};
  \node[st=yellow!20, minimum width=27mm] (synth) at (10.1,0) {\texttt{synth\_bus\_*}\\\texttt{synth\_tick}};
  \node[buf, minimum width=26mm] (cvl) at (14.5,0) {CV delay line\\\scriptsize $D+1$ entries};
  \node[st=green!15, minimum width=28mm] (tim1) at (18.8,0) {TIM1 table\\\scriptsize DMA ch2 $\to$ CV PWM};
  \node[buf, minimum width=26mm] (vq) at (6.6,-2.6) {$V_n$ queue\\\scriptsize stamped with $b_n$};
  \node[st=orange!18, minimum width=27mm] (rend) at (10.1,-2.6) {renderer\\\scriptsize DMA ch1 IRQ, low};
  \node[buf, minimum width=25mm] (dring) at (14.5,-2.6) {DAC ring\\\scriptsize $2\times256$ periods};
  \node[st=green!15, minimum width=28mm] (tim2) at (18.8,-2.6) {TIM2 CH2\\\scriptsize PWM DAC, PC2};
  \draw[->] (ring) -- (parse); \draw[->] (parse) -- (evq);
  \draw[->] (evq) -- node[above,lab]{due} (synth);
  \draw[->] (synth) -- node[above,lab]{\texttt{hw\_cv\_set}} (cvl);
  \draw[->] (cvl) -- node[above,lab]{tick $f+D$} (tim1);
  \draw[->] (synth.south) -- ++(0,-0.6) -| node[pos=0.25,above,lab]{\texttt{dsp\_commit}} (vq.north);
  \draw[->] (vq) -- node[below,lab]{at $b_n$} (rend);
  \draw[->] (rend) -- (dring);
  \draw[->] (dring) -- node[below,lab]{DMA} (tim2);
  \begin{scope}[on background layer]
    \node[draw=blue!50, dashed, rounded corners, fit=(parse)(evq)(synth), inner sep=2.5mm,
      label={[lab, blue!60!black]above:control tick of frame $f$ (high level)}] {};
  \end{scope}
  \node[lab, anchor=west, align=left] at (-1.1,-4.0) {bytes arrive during frame $f-1$ (sent at $f-L$) $\;\to\;$ run in
    the tick of frame $f$ $\;\to\;$ heard at frame $f+D$, on every card at the same carrier period};
\end{tikzpicture}}
\caption{Card data flow from a received event to the outputs. Buffers are grey; the dashed box runs in the control
tick, the renderer in the low-level DMA interrupt.}
\label{fig:sadflow}
\end{figure}

\paragraph{Bus state machine.} Figure~\ref{fig:sadstate} shows the card's states, as implemented in \file{bus.c}.
Each state's outputs: in STANDALONE the voice plays from the console or its tunes; in TRIM and LISTEN the receiver is
being prepared; in SYNCED timed events run; in LOST the VCA is muted until SYNC returns.

\begin{figure}[htbp]
\centering
\begin{tikzpicture}[font=\small\sffamily, >={Stealth[length=2mm]}, node distance=8mm and 20mm,
  s/.style={draw, rounded corners=6pt, fill=#1, minimum height=9mm, minimum width=24mm, align=center},
  s/.default=white, lab/.style={font=\scriptsize, align=center}]
  \node[s=gray!12] (sa) {STANDALONE\\\scriptsize SysTick ticks, RX off};
  \node[s=yellow!20, right=of sa] (tr) {TRIM\\\scriptsize 100 SYNC periods};
  \node[s=yellow!20, right=of tr] (li) {LISTEN\\\scriptsize RX on};
  \node[s=green!18, below=16mm of li] (sy) {SYNCED\\\scriptsize timed events run};
  \node[s=red!12, left=34mm of sy] (lo) {LOST\\\scriptsize VCA muted};
  \draw[->] ([xshift=-10mm]sa.west) node[left,lab]{reset} -- (sa);
  \draw[->] (sa) -- node[above,lab]{SYNC edge} (tr);
  \draw[->] (tr) to[loop above] node[above,lab]{HSITRIM step, measure again; a gap restarts} (tr);
  \draw[->] (tr) -- node[above,lab]{error $<\frac12$ step} (li);
  \draw[->] (li) -- node[right,lab]{first FRAME} (sy);
  \draw[->] ([yshift=2mm]sy.west) -- node[above,lab]{10 ticks without SYNC} ([yshift=2mm]lo.east);
  \draw[->] ([yshift=-2mm]lo.east) -- node[below,lab]{SYNC edge} ([yshift=-2mm]sy.west);
\end{tikzpicture}
\caption{Card bus states (\file{bus.c}). A missing SYNC edge is bridged by a SysTick tick and counted (it
does not leave SYNCED); an NRST or SWIO reset restarts at STANDALONE.}
\label{fig:sadstate}
\end{figure}

\paragraph{Memory.} Table~\ref{tab:sadmem} gives the card build's memory (link map of the current build). The DAC
ring, the renderer's work area and the console's capture buffer share one buffer (\texttt{hw\_scratch}): captures are
refused while the DAC runs.

\begin{table}[htbp]
\centering\small
\caption{Card build memory (CH32V006, \texttt{make BOARD=card}).}\label{tab:sadmem}
\begin{tabularx}{\textwidth}{@{}l r X@{}}
\toprule
Item & Bytes & Content \\
\midrule
flash: code and constants & 53\,228 & of 63\,488; plus the \texttt{.data} image: \SI{55.6}{KB} in all \\
RAM \texttt{.data} & 2\,416 & initialised data and about \SI{2.3}{KB} of render code run from SRAM \\
RAM \texttt{.bss} & 4\,608 & \texttt{hw\_scratch} 2560 (DAC ring 512, render work area 256, console capture);
  modal voices 760; RX ring 256; event queue 128; $V_n$ queue 128; CV delay line 78; packet and reply buffers 128;
  oscillator state 72; bus, patch and CV burst state 118; the rest small \\
stack region & 1\,168 & at least 656 bytes never used under bus load and stress, about 340 with the console \\
\bottomrule
\end{tabularx}
\end{table}

% ---------------------------------------------------------------------------------------------------------------------
\subsection{Controller firmware}\label{sad:ctrl}

\paragraph{Modules.} Table~\ref{tab:sadctmod} lists the controller's target modules (from \file{docs/controller.md})
and what the bench stand-in already provides.

\begingroup\small\setlength{\tabcolsep}{4pt}
\begin{longtable}{@{}>{\raggedright}p{3.0cm} >{\raggedright}p{7.6cm} >{\raggedright\arraybackslash}p{4.9cm}@{}}
\caption{Controller firmware modules: target and bench status.}\label{tab:sadctmod}\\
\toprule
Module & Responsibility & In \file{bench\_ctrl} \\
\midrule\endfirsthead
\toprule
Module & Responsibility & In \file{bench\_ctrl} \\
\midrule\endhead
\bottomrule\endfoot
frame engine & TIM1 CH1 SYNC; at each edge FRAME, enumeration, queued packets up to 95 bytes, refreshes; one TX DMA
  transfer & built \\
bus scheduler & stamps events with $a + L$; two queues (notes, then controllers); coalescing; late accounting &
  built (no coalescing) \\
reply handler & RX DMA ring on RETURN; COBS / CRC; HELLO, STATUS, COUNTERS, CAL\_DATA; time-outs & built \\
slot manager & 595 reset and release, enumeration, chip ID $\to$ slot, supervision and recovery (NRST, SWIO) & one
  slot, without NRST \\
USB device & composite USB-MIDI + CDC: USBD driver (backplane), USBFS driver (development board) & built (USBD
  enumerated by the own host) \\
USB host & class-compliant MIDI on USBFS, power switch and fault & built (\texttt{USBH=1}, \file{usbhost.c};
  self-loop tested, no hubs) \\
MIDI input & per port: running status, real time, SysEx assembly; ports: USB-MIDI, USB host, USART1 (DIN option) &
  USB-MIDI and the bench command \\
part router & channel and key range $\to$ part; CC, NRPN and sustain handling; program change; MIDI clock & one
  poly part, or one part per channel \\
allocator & voice choice within a slot mask; PATCH\_SELECT and PATCH\_DATA before a NOTE\_ON; unison and kit modes &
  basic allocator \\
patch store & 64 programs and setups in flash pages with wear levelling; working copies in SRAM & -- \\
calibration store & CAL\_DATA by chip ID & -- \\
CDC commands & text commands of midi.md 6; SysEx dumps & binary bench command set, identity reply \\
SWIO programmer & WCH one-wire protocol (bit-banged, \SI{125}{ns} resolution) on 8 lines; image from CDC, program and
  verify & Ardulink bridge for minichlink (one line) \\
supervision & watchdog; VBUS\_DET and the D+ pull-up; activity LED & -- \\
options & UAC1 audio from the ADC taps; CLK\_OUT on TIM3 CH4 & -- \\
\end{longtable}
\endgroup

\paragraph{Execution contexts.} The frame interrupt (TIM1 update, at the SYNC rising edge) has the highest priority
and builds each frame's transmission from queues filled by the main loop. Its budget is \SI{40}{\micro s} from the
edge to the start of the DMA; measured 0--3\,\textmu s with and without USB traffic. The USB interrupt has a lower
priority, and interrupt nesting is off, so USB can only delay the frame interrupt by its own length (at most
\SI{3}{\micro s} measured). The main loop serves the command ports, turns USB-MIDI packets into events, runs the
allocator and, in the target design, the storage and programming tasks. A flash erase or a SWIO programming run must
not hold up the queues: both run in the main loop between frames, and the frame interrupt keeps sending refreshes.

\begin{figure}[htbp]
\centering
\resizebox{\textwidth}{!}{%
\begin{tikzpicture}[font=\small\sffamily, >={Stealth[length=2mm]},
  st/.style={draw, rounded corners=2pt, fill=#1, minimum height=10mm, minimum width=24mm, align=center},
  st/.default=white, buf/.style={draw, fill=gray!10, minimum height=9mm, minimum width=21mm, align=center},
  lab/.style={font=\scriptsize}]
  \node[st=gray!12] (usbd) at (0,1.5) {USB-MIDI\\\scriptsize USBD};
  \node[st=gray!12] (host) at (0,0) {USB host\\\scriptsize keyboard};
  \node[st=gray!12] (cdc) at (0,-1.5) {CDC, SysEx\\\scriptsize editor};
  \node[st=blue!12] (midi) at (3.6,0) {MIDI input\\\scriptsize per port};
  \node[st=blue!12, minimum width=27mm] (parts) at (6.9,0) {part router\\\scriptsize CC, NRPN, sustain};
  \node[st=blue!12] (alloc) at (10.2,0) {allocator\\\scriptsize slot masks};
  \node[buf] (q1) at (13.4,1.0) {queue 1\\\scriptsize notes};
  \node[buf] (q2) at (13.4,-1.0) {queue 2\\\scriptsize controllers};
  \node[st=orange!18, minimum width=26mm] (fe) at (17.0,0) {frame engine\\\scriptsize TIM1 IRQ, 1\,kHz};
  \node[st=green!15, minimum width=26mm] (tx) at (20.6,0) {USART3 TX\\\scriptsize DMA $\to$ BCAST};
  \node[st=blue!12, minimum width=27mm] (swio) at (3.6,-3.3) {SWIO programmer\\\scriptsize 8 lines};
  \node[st=blue!12, minimum width=27mm] (store) at (6.9,-3.3) {patch store\\\scriptsize flash};
  \node[st=blue!12, minimum width=27mm] (slot) at (13.4,-3.3) {slot manager\\\scriptsize 595, HELLO};
  \node[st=green!15, minimum width=26mm] (rx) at (20.6,-3.3) {USART3 RX\\\scriptsize RETURN};
  \draw[->] (usbd.east) -- ++(0.5,0) |- ([yshift=2mm]midi.west);
  \draw[->] (host.east) -- (midi.west);
  \draw[->] (cdc.east) -- ++(0.5,0) |- ([yshift=-2mm]midi.west);
  \draw[->] (midi) -- (parts); \draw[->] (parts) -- (alloc);
  \draw[->] (alloc.east) -- ++(0.5,0) |- (q1.west);
  \draw[->] (alloc.east) -- ++(0.5,0) |- (q2.west);
  \draw[->] (q1.east) -- ++(0.6,0) |- ([yshift=2mm]fe.west);
  \draw[->] (q2.east) -- ++(0.6,0) |- ([yshift=-2mm]fe.west);
  \draw[->] (fe) -- (tx);
  \draw[->] (parts.south) -- node[right,lab]{programs} (store.north);
  \draw[->] (rx.west) -- node[below,lab]{replies} (slot.east);
  \draw[->] (slot.north) -- ++(0,0.6) -| node[pos=0.25,above,lab]{enumeration, refreshes} (fe.south);
  \draw[->] (cdc.south) |- node[pos=0.75,below,lab]{card image} (swio.west);
\end{tikzpicture}}
\caption{Controller data flow (target). Blue modules run in the main loop, the frame engine in the frame interrupt.}
\label{fig:sadctflow}
\end{figure}

\paragraph{Start-up sequence.} Clocks and USB; shift 0x00 into the 595 and enable its outputs (all cards in reset);
start SYNC and FRAME; enumerate slot by slot (bus.md 6); send CONFIG and the active patches; then start USB-MIDI
processing. A controller reset at any point restarts this sequence; the cards run muted in the meantime (A10).

\paragraph{Storage (proposed).} Flash pages for 64 programs of 144 bytes (\SI{9.2}{KB}), 8 setups and the
calibration copies (8 cards $\times$ about 64 bytes), written by appending to a log with a sequence number and a CRC
and compacted when full. The working copies in SRAM are what the parts use. A corrupt or empty store falls back to
built-in defaults.

% ---------------------------------------------------------------------------------------------------------------------
\subsection{Host software}\label{sec:sadhs}

The browser app (requirements in Section~\ref{sec:srshs}) is a set of static files in the repository (\file{app/};
first version 2026-10-20): plain HTML and JavaScript without a build step and without modules (Chromium blocks
module scripts from \texttt{file://}), so it opens from disk or from any static web server and works offline (A12).
An emulated instrument (programs, parts, six voice cards, voice allocation, Web Audio sound) stands in for the
hardware; the tests run in Node (protocol and emulator) and in headless Edge (an offline render of the sound). Its
layers:

\begin{table}[htbp]
\centering\small
\caption{Host software modules (\file{app/js/}).}\label{tab:sadhs}
\begin{tabularx}{\textwidth}{@{}l X@{}}
\toprule
Module & Responsibility \\
\midrule
MIDI link & Web MIDI: find the instrument's ports, send NRPN and SysEx, receive replies and incoming notes; SysEx
  framing and checksums \\
serial link & Web Serial on the CDC port: the maintenance commands, firmware image transfer, the console bridge \\
instrument model & the state mirrored from the instrument: active patch, programs, parts, cards and their status;
  every view reads and changes it, the links keep it in sync \\
parameter table & generated from the shared JSON (A13): IDs, names, ranges, units, CC numbers, grouping for the
  editor \\
views & patch editor, librarian, setup (parts), instrument view, maintenance, keyboard and MIDI monitor \\
files & import and export of \texttt{.syx} patches and banks, firmware images chosen by the user \\
\bottomrule
\end{tabularx}
\end{table}

\paragraph{Protocol split.} Everything musical goes over MIDI, so a DAW or any MIDI host reaches the same functions
(CT13): live parameter changes as NRPN (midi.md 3), patches, banks, setups, parts, card status and calibration as
SysEx messages with the non-commercial ID 0x7D, each with a request and a reply. Maintenance goes over the CDC port:
firmware images for the cards (programmed by the controller over SWIO) and for the controller, the console bridge
to one card, bulk diagnostics. The SysEx message set is drafted in \file{docs/sysex.md} (implemented by the
emulator); the maintenance commands are to be specified with the controller firmware.

\paragraph{Development path.} The app can be written and tested against the bench controller before the backplane
exists: it already is a USB-MIDI and CDC device, and the browser on the lab PC reaches both. The SysEx and NRPN
handling then moves with the controller firmware from the bench to the backplane unchanged.

% ---------------------------------------------------------------------------------------------------------------------
\subsection{Interfaces}

\begin{table}[htbp]
\centering\small
\caption{Software interfaces.}\label{tab:sadif}
\begin{tabularx}{\textwidth}{@{}l l X@{}}
\toprule
Interface & Between & Specification and state \\
\midrule
voice bus & controller, cards & \file{docs/bus.md} 2--6 (0.2, implemented; card firmware version 3) \\
patch parameter IDs & controller, cards, editor & \file{docs/midi.md} 4 (proposal; the cards implement a subset) \\
USB-MIDI & computer, controller & class compliant, one cable each way; \file{docs/controller.md} 5. USB ID 1209:0001
  is a test ID for development only \\
CDC & computer, controller & bench: binary commands (\file{bench\_ctrl.c} header); target: text commands
  (\file{docs/midi.md} 6, to define) \\
SysEx & computer, controller & ID 0x7D: identity reply implemented; patches, banks, setups, parts, card status
  and calibration to define (Section~\ref{sec:sadhs}) \\
NRPN & computer, controller & live parameter changes, IDs of midi.md 4 (from the shared parameter table, A13) \\
Web MIDI, Web Serial & browser app, computer & the app's access to the instrument's MIDI and CDC ports \\
SWIO console & computer, card & \file{console.c} commands (\file{firmware/voice/README.md}) \\
Ardulink & computer, controller & minichlink's Ardulink programmer protocol, for card programming through the
  controller \\
generated tables & tools, card firmware & \file{tables.h} from the result files (\file{gen\_tables.py}) \\
\bottomrule
\end{tabularx}
\end{table}

% ---------------------------------------------------------------------------------------------------------------------
\subsection{Build variants}

\begin{table}[htbp]
\centering\small
\caption{Build variants.}\label{tab:sadbuild}
\begin{tabularx}{\textwidth}{@{}l l X@{}}
\toprule
Firmware & Option & Effect \\
\midrule
card & \texttt{BOARD=card} & the voice card's pin map and the bus receiver (default: board~1 with its DCO and MIDI
  input) \\
 & \texttt{BUS\_TX\_PP=1} & bench only: the return line driven push-pull during a reply (no \SI{1}{k\ohm} pull-up on
  the bench) \\
 & \texttt{BUS\_BAUD=2000000} & \SI{2}{Mbaud} bus (both firmwares must match) \\
 & \texttt{DAC\_DUAL=1} & dual-PWM DAC (high and low byte, PC1) \\
 & chip & CH32V006 (card, board~1 with a V006); V003 and V002 builds compile \\
controller (bench) & \texttt{USB=usbd} & USB on USBD, PA11 / PA12 (the backplane's port; default USBFS on the
  development board) \\
 & \texttt{BUS\_BAUD=2000000} & as above \\
\bottomrule
\end{tabularx}
\end{table}

% ---------------------------------------------------------------------------------------------------------------------
\subsection{Verification architecture}

The software is verified at five levels (Table~\ref{tab:sadver}). Each requirement of Section~\ref{sec:srs} names
the level that covers it; Table~\ref{tab:srssum} summarises their status.

\begin{table}[htbp]
\centering\small
\caption{Verification levels.}\label{tab:sadver}
\begin{tabularx}{\textwidth}{@{}l l X@{}}
\toprule
Level & Tool & Covers \\
\midrule
model & \file{sim/behav/} (\file{dsp.py}, \file{behav.py}) & algorithms and sound, against the analog model (S9) \\
host regression & \file{firmware/voice/test/render\_test.py} & the renderer bit-exact against a reference revision,
  197 scenarios \\
on target & console \texttt{check}, \texttt{bench} & assembly and compiler (hash), cycle counts per patch \\
bench integration & \file{tools/bustest.py}, \file{tools/usbtest.py} & protocol, timing, latency, errors, load with 7
  virtual voices, USB-MIDI and CDC; V203 development board and one V006 \\
system & stages V2--V4 (Table~\ref{tab:vstages}) & real cards and backplane: NRST, alignment between cards,
  calibration, 8 cards \\
\bottomrule
\end{tabularx}
\end{table}

\begin{table}[htbp]
\centering\small
\caption{Status of the software requirements (Section~\ref{sec:srs}), 20 October 2026.}\label{tab:srssum}
\begin{tabular}{@{}l r r r r r@{}}
\toprule
Group & Requirements & met & partial & open & design \\
\midrule
card firmware (CF) & 23 & 16 & 4 & 3 & 0 \\
voice bus (BUS) & 8 & 5 & 1 & 0 & 2 \\
controller firmware (CT) & 22 & 2 & 11 & 8 & 1 \\
tools, non-functional (TL, NF) & 8 & 8 & 0 & 0 & 0 \\
host software (HS) & 11 & 1 & 9 & 1 & 0 \\
\midrule
all & 72 & 32 & 25 & 12 & 3 \\
\bottomrule
\end{tabular}
\end{table}

% ---------------------------------------------------------------------------------------------------------------------
\subsection{Open architecture questions}\label{sec:sadopen}

\begin{enumerate}
  \item \textbf{USB stack for the host port.} Extend the own drivers (A8) to a minimal MIDI host, or port TinyUSB
    for both ports? The host port is fitted now, and a minimal MIDI host on the own drivers works in a self-loop on
    the development board (Section~\ref{sec:fw}); TinyUSB would add hubs and more device classes.
  \item \textbf{Parsing location (A4).} Stay with parsing in the tick ($L = 2$), or move it to the main loop now
    ($L = 3$, one more millisecond of latency) to decouple the bus load from the render for good?
  \item \textbf{CDC protocol.} Musical functions go over MIDI and maintenance over CDC (decided, 2026-10-20). Open:
    the maintenance commands as text or binary, and whether the bench's binary command set stays for the tests.
  \item \textbf{USB identity.} 1209:0001 is the pid.codes test ID. A product ID (pid.codes or another VID) is
    needed before the instrument leaves the bench.
  \item \textbf{Controller self-update.} The factory bootloader through BOOT0 (untested on the USB-C port), or an
    update mode in the firmware over CDC with a small resident loader?
  \item \textbf{Card firmware compatibility.} With cards updated one by one, should the controller refuse to start
    cards with a different firmware version (HELLO) or update them automatically?
\end{enumerate}
